% =====================================================================
%  Seance 5 -- D. fiabilite du detecteur + limites (obligatoire)
% =====================================================================
\section{Fiabilité du détecteur et limites}

\begin{frame}{Fiabilité 1 : la dérive de modèle}
\footnotesize%
\begin{columns}[t]
\begin{column}{0.52\textwidth}
\begin{block}{Le trafic change, le modèle vieillit}
\begin{itemize}
\item Lundi : nouveau service batch qui se connecte toutes les nuits à 3 h.
\item Un mois plus tard : rotation des mots de passe, politiques de comptes
modifiées, croissance du parc.
\item Le « normal » appris en octobre n'est plus le normal de décembre :
les scores dérivent, les faux positifs s'accumulent \textbf{sans aucune
erreur visible}.
\end{itemize}
\end{block}
\begin{ebox}[Les garde-fous]
\begin{itemize}
\item Surveiller la \textbf{distribution des features} elles-mêmes
(PSI, KS test) : alarme si le normal bouge.
\item Réentraîner sur fenêtre glissante (30-90 j) --- et garder une
référence figée pour comparer.
\item Chaque incident confirmé $=$ étiquette neuve (boucle du cours 2).
\end{itemize}
\end{ebox}
\end{column}
\begin{column}{0.44\textwidth}
\begin{kbox}[Lien avec les séances précédentes]
C'est la \textbf{même dérive} que le capteur de srv-06 (séance 2), mais sur
les features de sécurité au lieu de la température. La leçon est la même :
\textbf{un modèle non surveillé est un capteur qui ment paisiblement}.
\end{kbox}
\begin{wbox}[Le piège du réentraînement]
Si l'attaquant revient toutes les semaines, le réentraînement glissant
\textbf{apprend l'attaque comme normal}. Fixer la fenêtre, surveiller les
scores, garder l'humain dans la boucle.
\end{wbox}
\end{column}
\end{columns}
\end{frame}

\begin{frame}{Fiabilité 2 : la qualité des données en production}
\footnotesize%
\begin{columns}[t]
\begin{column}{0.5\textwidth}
\begin{block}{Ce qui casse un détecteur de sécurité sans prévenir}
\begin{itemize}
\item \textbf{Horloges} : un serveur en retard déplace ses événements de
fenêtre --- fausse rafale, vrai faux positif.
\item \textbf{Collecte perdue} : agent tombé $=$ fenêtres vides $=$
aveugle pendant l'attaque. Surveiller le \emph{volume} de logs par machine !
\item \textbf{Champ trompeur} : dans nos logs, des lignes « Rotation des
logs » ou « Sauvegarde » portent aussi \py{auth\_reussie = False}.
Compter bêtement les \py{False} donne 4\,828 échecs dont beaucoup ne sont
pas des authentifications : \textbf{filtrer d'abord, compter ensuite}.
\end{itemize}
\end{block}
\end{column}
\begin{column}{0.46\textwidth}
\begin{block}{Les réflexes de la séance 2, rejoués en sécurité}
\begin{itemize}
\item Compte de lignes par machine et par jour (trous de collecte).
\item Validation de schéma (types, niveaux attendus).
\item Vérification croisée : le volume d'échecs d'une machine ne peut pas
être 10$\times$ la médiane sans que \emph{quelque chose} se passe ---
mais est-ce une attaque ou un défaut d'agent ?
\end{itemize}
\end{block}
\begin{kbox}[Règle]
\textbf{Un détecteur est aussi fiable que sa pipeline de données.}
Le ML ne corrige pas un capteur menteur.
\end{kbox}
\end{column}
\end{columns}
\end{frame}

\begin{frame}{Fiabilité 3 : la redondance des détecteurs}
\begin{center}
\footnotesize
\begin{tabular}{@{}p{3.6cm}p{3.4cm}p{3.4cm}@{}}
\toprule
 & \textbf{Règle (seuil 5 échecs/10 min)} & \textbf{Modèle (7 features)}\\
\midrule
Attaque grossière (notre brute force) & \ok{} détecte (12/12) & \ok{} détecte (12/12)\\
\addlinespace[1pt]
Attaque lente (4 échecs/10 min) & \nok{} invisible & \ok{} visible (entropie, comptes, nuit)\\
\addlinespace[1pt]
Bruit opérationnel (défauts de service) & \nok{} faux positifs possibles & \ok{} absorbé par le contexte\\
\addlinespace[1pt]
Explication à donner & \ok{} triviale & $\approx$ SHAP (slide suivante)\\
\addlinespace[1pt]
Coût de maintenance & règle à réviser à chaque changement & réentraînement, surveillance de dérive\\
\bottomrule
\end{tabular}
\end{center}
\vspace{1mm}
\begin{kbox}[Défense en profondeur, version ML]
Règles et modèle fonctionnent \textbf{en parallèle}, jamais l'un à la place
de l'autre : la règle couvre le cas trivial et sert de référence explicite ;
le modèle couvre l'espace que la règle ne sait pas formuler. Une alerte
confirmée par les deux canaux monte en priorité.
\end{kbox}
\end{frame}

\begin{frame}{Limite 1 : l'attaque adversariale}
\begin{columns}[t]
\begin{column}{0.52\textwidth}
\begin{block}{L'attaquant connaît (ou teste) votre détecteur}
S'il sait que le top du classement regarde échecs, rafales et comptes
sensibles, il s'adapte :
\begin{itemize}
\item ralentir : 2 échecs / 10 min, attaque étalée sur 20 h ;
\item doper l'entropie : tester beaucoup de comptes plausibles, même faux ;
\item passer par des comptes ordinaires (\py{dev1}) au lieu de \py{root}.
\end{itemize}
Chaque contremesure dégrade \emph{les mêmes features} que le modèle utilise.
\end{block}
\end{column}
\begin{column}{0.44\textwidth}
\begin{ebox}[Ce que ça coûte à l'attaquant]
Ralentir à 2 échecs/10 min multiplie la durée de l'attaque par 5
(2 h $\rightarrow$ 10 h pour 125 essais) : plus de temps de détection
par d'autres canaux, plus de chances qu'un gardien remarque les connexions
de nuit. \textbf{Le ML ne stoppe pas l'attaquant, il le coûteux.}
\end{ebox}
\begin{kbox}[La parade]
Horloge d'émission, cadence inter-essais, source réseau, corrélations entre
machines. Varier et \textbf{ne pas documenter} toutes les dimensions de détection.
\end{kbox}
\end{column}
\end{columns}
\end{frame}

\begin{frame}{Limite 2 : la fatigue d'alerte --- Limite 3 : l'explicabilité}
\footnotesize%
\begin{columns}[t]
\begin{column}{0.5\textwidth}
\begin{block}{Fatigue d'alerte}
Trop de faux positifs $\Rightarrow$ les analystes \textbf{clichent sans
regarder} $\Rightarrow$ la vraie attaque passe au milieu. C'est ce qui s'est
passé à Target (2013) : les alarmes existaient, personne ne les lisait plus.
\begin{itemize}
\item Mesure : \textbf{taux de traitement} des alertes, précision@N
\emph{en conditions réelles}.
\item Remède : classer (précision@N), regrouper (clustering), seuil haut +
revue périodique du rang 30-100.
\end{itemize}
\end{block}
\end{column}
\begin{column}{0.46\textwidth}
\begin{block}{Explicabilité : SHAP en un aperçu}
SHAP attribue à \textbf{chaque feature sa part} dans la décision d'une
session donnée. Pour notre session n°1 (15 échecs, tous sur comptes
sensibles) : l'attribution pointe d'abord les échecs sur comptes sensibles,
puis la rafale et le taux d'échec --- cohérent avec les importances
mesurées à la slide précédente.
L'analyste reçoit non pas « anormal » mais \emph{pourquoi}. Un modèle
explicable est un modèle qu'on ose couper en incident.
\end{block}
\end{column}
\end{columns}
\vspace{1mm}
\begin{wbox}[Le cercle vicieux à éviter]
Alertes non traitées $\rightarrow$ étiquettes jamais produites $\rightarrow$
détecteur supervisé impossible à améliorer $\rightarrow$ encore plus de faux
positifs. Traiter les alertes, c'est \textbf{alimenter} le système.
\end{wbox}
\end{frame}

\begin{frame}{Limite 4 : la vie privée des logs}
\begin{columns}[t]
\begin{column}{0.5\textwidth}
\begin{block}{Un log d'authentification est une donnée personnelle}
Il dit \emph{qui} était où, quand, sur quelle machine. Nos colonnes :
\py{utilisateur}, \py{machine}, \py{horodatage} --- la combinaison reconstitue
des habitudes de travail (horaires, localisation, usages).
\end{block}
\begin{itemize}
\item \textbf{Pseudonymisation} : remplacer \py{dev1} par un identifiant
op$\;$\#417 (table de correspondance hors zone d'analyse, accès restreint).
\item Les features calculées (échecs, entropie, compte \emph{nombre} de
comptes) survivent à la pseudonymisation : le ML n'a pas besoin des noms.
\item Durée de conservation minimale, accès journalisés, minimisation des
colonnes.
\end{itemize}
\end{column}
\begin{column}{0.46\textwidth}
\begin{ebox}[Notre TP, version conforme]
\begin{itemize}
\item Entraîner le détecteur sur les \textbf{features}, pas sur les noms.
\item Ne révéler les identités que dans la \textbf{note d'incident}, après
corrélation autorisée --- c'est le livrable B5.
\item Les comptes \py{inconnu-1}, \py{inconnu-2} du jeu sont déjà l'exemple
idéal : des pseudos, et l'attaque reste parfaitement détectable.
\end{itemize}
\end{ebox}
\begin{kbox}[Éthique et efficacité vont ensemble]
Un détecteur pseudonymisé est \emph{moins} exposé aux abus internes ---
et fonctionne aussi bien. Penser privacy by design n'est pas une contrainte
de plus : c'est une architecture.
\end{kbox}
\end{column}
\end{columns}
\end{frame}
